\section{Usage and Examples}
This section provides a minimal example that emits a signal and awaits its
presence in the next instant, followed by a larger example that uses guarded
execution and weak preemption. The examples are aligned with the tests in
\texttt{tests/ok} to ensure they are executable and representative (e.g.,
\texttt{emit\_once\_await\_one.ml}, \texttt{watch\_when\_complete.ml}).

To execute the examples, run the corresponding test binaries, e.g.:
\begin{lstlisting}[style=tempo]
dune exec tests/ok/emit_once_await_one.exe
dune exec tests/ok/watch_when_complete.exe
\end{lstlisting}

Reading the examples. Each example is written as a function \texttt{program}
passed to \texttt{execute}. The host provides an input stream by emitting on
the input signal before each instant and observes output after each instant.
Thus, the observable behavior of a program is a sequence of values emitted on
the output signal at instant boundaries.

\paragraph{End-to-end example 1: input/output echo}
This example waits for an input value and echoes it on the output signal in the
next instant.
\begin{lstlisting}[style=tempo]
let program input output =
  let v = await input in
  emit output v

let () = execute program
\end{lstlisting}
Step-by-step behavior:
\begin{enumerate}
\item Instant $t$: the host emits on \texttt{input} via the \texttt{execute}
  hook; the program is awaiting and either registers a waiter or is scheduled
  for the next instant if the signal is already present.
\item Instant $t{+}1$: the awaiting continuation resumes and emits on
  \texttt{output}, which the host observes at the end of the step.
\end{enumerate}

\paragraph{End-to-end example 2: guarded wait with timeout}
This example waits for a request, but is weakly preempted if a timeout signal
is emitted after $n$ instants.
\begin{lstlisting}[style=tempo]
let rec timeout_after n s =
  if n = 0 then emit s ()
  else (pause (); timeout_after (n - 1) s)

let program input output =
  let request = new_signal () in
  let timeout = new_signal () in
  let _ = fork (fun () -> timeout_after 3 timeout) in
  let handled = new_signal () in
  let _ = fork (fun () ->
    watch timeout (fun () ->
      let v = await request in
      emit handled ();
      emit output (`Ok v))) in
  let _ = fork (fun () ->
    when_ timeout (fun () ->
      present_then_else handled
        (fun () -> ())
        (fun () -> emit output `Timeout))) in
  let v = await input in
  emit request v

let () = execute program
\end{lstlisting}
Step-by-step behavior:
\begin{enumerate}
\item A timer process emits \texttt{timeout} after 3 instants.
\item A forked branch waits on \texttt{request} but is wrapped in
  \texttt{watch timeout}, so it is terminated at the end of the instant in
  which \texttt{timeout} is emitted.
\item If \texttt{request} arrives before timeout, the guarded branch emits
  \texttt{\`Ok v} and marks \texttt{handled}. Otherwise, the timeout branch
  emits \texttt{\`Timeout}.
\end{enumerate}
Note: the \texttt{handled} guard prevents emitting \texttt{\`Timeout} in the
same instant as \texttt{\`Ok v}, even under weak preemption.

In practice, this pattern makes timeouts compositional: the timeout branch is
expressed separately from the request handler, while determinism is preserved
through explicit guards and instant boundaries.

\paragraph{Running example used in the implementation section}\label{ex:guard-watch}
The following small program appears again in the implementation section to
illustrate queue movement, guard blocking, and weak preemption:
\begin{lstlisting}[style=tempo]
let g = new_signal () in
let s = new_signal () in
let guarded = (fun () -> when_ g (fun () -> emit s ())) in
let watched = (fun () -> watch s (fun () -> pause (); emit s ())) in
parallel [guarded; watched]
\end{lstlisting}
The key behaviors are: (1) the guarded task blocks until \texttt{g} is present;
(2) the watched task can emit \texttt{s} in the same instant; (3) the
\texttt{watch} prevents the watched task from continuing in the next instant.

\paragraph{Guard disjunctions and preemption}
The following patterns illustrate that guard disjunctions are obtained by
parallel guarded branches, and that nested preemptions act as a disjunction of
stopping conditions.
\begin{lstlisting}[style=tempo]
(* Guard disjunction via parallel branches. *)
let guarded body =
  parallel [
    (fun () -> when_ g1 (fun () -> body ()));
    (fun () -> when_ g2 (fun () -> body ()))
  ]

(* Preemption disjunction: stop if s1 OR s2 is emitted. *)
let preempted body =
  watch s1 (fun () ->
    watch s2 (fun () ->
      body ()))
\end{lstlisting}

\paragraph{Derived control primitives}
\tempo{} provides higher-level control structures built from the primitives.
For instance, \linebreak \texttt{present\_then\_else} branches on signal
presence at the end of an instant. It provides a deterministic conditional
construct.
\begin{lstlisting}[style=tempo]
(* Execute then_body if s becomes present this instant, else else_body. *)
present_then_else s
  (fun () -> then_body ())
  (fun () -> else_body ())
\end{lstlisting}

Another derived pattern is \texttt{every\_do}, which executes a body whenever
the signal becomes present, while avoiding immediate re-triggering within the
same instant.
\begin{lstlisting}[style=tempo]
let rec every_do s body =
  let _ = await s in
  pause ();
  let completion = new_signal () in
  watch s (fun () ->
    body ();
    emit completion ());
  present_then_else completion
    (fun () -> pause (); every_do s body)
    (fun () -> every_do s body)
\end{lstlisting}
This pattern uses a \texttt{pause} to avoid strong re-triggering in the same
instant, preserving weak restart semantics aligned with synchronous control.
This corresponds to the restart discipline described in Section~4, where
explicit pauses separate instants to keep reactions deterministic.

\paragraph{Use case: reactive I/O protocol}
This example models a simple device handshake: a request from the environment
must be acknowledged within a bounded number of instants. If the acknowledgement
does not arrive in time, the protocol emits a timeout. This pattern is typical
of sensor/actuator control loops where a response deadline is required.
Here, the environment is the host input stream provided to \texttt{execute}.
The device model is a simulated internal component used to illustrate the
protocol; in a real system it would be replaced by actual I/O.
\begin{lstlisting}[style=tempo]
let rec timeout_after n timeout =
  if n = 0 then emit timeout ()
  else (pause (); timeout_after (n - 1) timeout)

let protocol request ack output =
  let timeout = new_signal () in
  let _ = fork (fun () -> timeout_after 2 timeout) in
  let handled = new_signal () in
  let _ = fork (fun () ->
    watch timeout (fun () ->
      let payload = await request in
      emit handled ();
      emit output (`Got payload))) in
  let _ = fork (fun () ->
    when_ timeout (fun () ->
      present_then_else handled
        (fun () -> ())
        (fun () -> emit output `Timeout))) in
  let _ = await ack in
  ()

let program input output =
  let ack = new_signal () in (* internal ack from device model *)
  let _ = fork (fun () -> protocol input ack output) in
  (* Environment-provided request; internal device emits ack. *)
  let _ = fork (fun () ->
    let _ = await input in
    pause ();
    emit ack ()) in
  ()

let () = execute program
\end{lstlisting}
Step-by-step behavior:
\begin{enumerate}
\item Instant $t$: the protocol waits for \texttt{request} and arms a timer.
\item Instant $t{+}1$: if \texttt{request} arrives, the simulated device emits
  \texttt{ack} in the next instant; the protocol can emit \texttt{`Got payload}
  unless the timeout fires first.
\item Instant $t{+}2$: if no request was handled, the timeout branch emits
  \texttt{`Timeout}. The \texttt{watch} ensures the request handler is
  preempted once the timeout is present.
\end{enumerate}

\paragraph{Timeline (environment vs internal)}
\begin{table}[h]
\centering
\begin{tabularx}{\linewidth}{lY}
\hline
Instant & Observable events \\ \hline
$t$ & Environment may emit \texttt{request}. \\
$t{+}1$ & Internal: \texttt{ack} emitted by device model; protocol may emit
\texttt{`Got payload}. \\
$t{+}2$ & Internal: timeout may emit \texttt{`Timeout} if not handled. \\ \hline
\end{tabularx}
\caption{Timeline for the reactive I/O protocol.}
\end{table}

\paragraph{Extended scenario: periodic sampling with timeout}
This example models a periodic sensor read where a sample must be acknowledged
within a fixed number of instants; otherwise a timeout is raised.
\begin{lstlisting}[style=tempo]
let rec every_tick tick body =
  let _ = await tick in
  body ();
  pause ();
  every_tick tick body

let rec timeout_after n timeout =
  if n = 0 then emit timeout ()
  else (pause (); timeout_after (n - 1) timeout)

let program input output =
  let tick = new_signal () in
  let ack = new_signal () in
  let timeout = new_signal () in
  let _ = fork (fun () -> every_tick tick (fun () -> emit input ())) in
  let _ = fork (fun () -> timeout_after 2 timeout) in
  let _ = fork (fun () ->
    watch timeout (fun () ->
      let _ = await input in
      emit ack ())) in
  when_ timeout (fun () -> emit output `Timeout);
  let _ = await ack in
  emit output `Ok
\end{lstlisting}
This scenario highlights the interaction between periodic activation, weak
preemption, and bounded response times.
